\section{Proofs}
\label{app:identifiability_proofs}

In this appendix \(A\) denotes the projected operator \(\widetilde A\) of
Equation~\ref{eq:projected_model}, \(\mathbf a_j\) its columns,
\(V_g=\operatorname{col}(A_g)\), \(W_g=\operatorname{col}(A_{E\setminus g})\) and
\(U=\operatorname{col}(A)\).  The unit tests in the code release
(\texttt{tests/test\_iasa\_grouping.py}) check numerically the exact components against
brute force (Theorem~\ref{thm:groupings}(b)--(d)), the two-column separation
\(\sqrt{1-\rho^2}\), Theorem~\ref{thm:group_error}(a), the examples of
Proposition~\ref{prop:reportable_structure}, and the exhaustive and greedy partition
searches.

\subsection{Individual Contributions}

\noindent\textbf{Proof of Proposition~\ref{prop:rank_identifiability}.}
If \(\operatorname{rank}(A)=J\), the rank--nullity theorem gives
\(\ker(A)=\{\mathbf 0\}\), so \(A\mathbf c_1=A\mathbf c_2\) implies
\(\mathbf c_1=\mathbf c_2\).  Conversely, if \(\operatorname{rank}(A)<J\) there is a
nonzero \(\mathbf d\in\ker(A)\).  Choose \(\beta>\|\mathbf d\|_\infty\) and set
\(\mathbf c=\beta\mathbf 1\).  Then \(\mathbf c\) and \(\mathbf c+\mathbf d\) are
distinct, both lie in \(\mathbb R_+^J\) because
\(\beta+d_j\ge\beta-\|\mathbf d\|_\infty>0\), and
\(A(\mathbf c+\mathbf d)=A\mathbf c\).  \(\square\)

The same interior-point construction is used below whenever a kernel vector must be
turned into two nonnegative vectors with equal projected observations.

\noindent\textbf{Proof of Proposition~\ref{prop:noise_robust}.}
For least squares with full column rank,
\(\widehat{\mathbf c}-\mathbf c=A^\dagger\widetilde{\boldsymbol\epsilon}
=A^\dagger P_A\widetilde{\boldsymbol\epsilon}\), and
\(\|A^\dagger\|_2=1/\sigma_J(A)\).  For NNLS, let
\(C=\{A\mathbf z:\mathbf z\ge0\}\), a closed convex cone contained in \(U\).  For
\(\mathbf x\in U\),
\(\|\yt-\mathbf x\|_2^2=\|P_A\yt-\mathbf x\|_2^2+\|(I-P_A)\yt\|_2^2\), so the
fitted signal is \(A\widehat{\mathbf c}=\Pi_C(\yt)=\Pi_C(P_A\yt)\), where \(\Pi_C\) is
the Euclidean projection onto \(C\).  Since \(\mathbf c\ge0\), \(A\mathbf c\in C\)
and \(\Pi_C(A\mathbf c)=A\mathbf c\).  Projection onto a closed convex set is
nonexpansive, so
\[
\|A(\widehat{\mathbf c}-\mathbf c)\|_2
=\|\Pi_C(P_A\yt)-\Pi_C(A\mathbf c)\|_2
\le\|P_A\yt-A\mathbf c\|_2
=\|P_A\widetilde{\boldsymbol\epsilon}\|_2 .
\]
Full column rank gives \(\|A\Delta\|_2\ge\sigma_J(A)\|\Delta\|_2\) for
\(\Delta=\widehat{\mathbf c}-\mathbf c\), which proves the first inequality
in~\eqref{eq:nnls_noise_bound}; the second holds because \(P_A\) is an orthogonal
projector.  \(\square\)

\noindent\textbf{Ray distance (Equation~\ref{eq:ray_distance_coherence_relation}).}
For nonzero \(\mathbf a_i,\mathbf a_j\), the function
\(\alpha\mapsto\|\mathbf a_i-\alpha\mathbf a_j\|_2^2/\|\mathbf a_i\|_2^2\) is
strictly convex with minimizer
\(\alpha^\star=\mathbf a_i^\top\mathbf a_j/\|\mathbf a_j\|_2^2\) and minimum
\(1-\rho_{ij}^2\).  The minimum is also \(\sin^2\) of the angle between
\(\mathbf a_i\) and \(\operatorname{span}(\mathbf a_j)\), so for \(J=2\) the ray
distance equals the sine of the angle between the two columns.  \(\square\)

\noindent\textbf{Proof of Proposition~\ref{prop:operator_perturbation}.}
Let \(\widehat A\) have full column rank and \(\|\widehat A-A^\star\|_2\le\delta_H\).
The data obey
\(\yt=A^\star\mathbf c+\widetilde{\boldsymbol\epsilon}
=\widehat A\mathbf c+[\widetilde{\boldsymbol\epsilon}+(A^\star-\widehat A)\mathbf c]\),
and \(\widehat A^\dagger\widehat A=I_J\), so
\(\widehat{\mathbf c}-\mathbf c=\widehat A^\dagger[\widetilde{\boldsymbol\epsilon}
+(A^\star-\widehat A)\mathbf c]\).  Taking norms gives
Equation~\ref{eq:operator_error_bound}.  \(\square\)

\subsection{Identifiable Groupings}

\noindent\textbf{Proof of Theorem~\ref{thm:groupings}.}
\emph{(a)}  Suppose \(U=\bigoplus_gV_g\).  If \(A\mathbf c_1=A\mathbf c_2\), then
\(\sum_gA_g(\mathbf c_{1,g}-\mathbf c_{2,g})=0\) with each term in \(V_g\), so each
term is zero.  Conversely, if the sum \(\sum_gV_g\) is not direct, there are
\(\mathbf u_g\in V_g\), not all zero, with \(\sum_g\mathbf u_g=0\).  Writing
\(\mathbf u_g=A_g\mathbf d_g\) defines \(\mathbf d\in\ker(A)\) with
\(A_h\mathbf d_h\ne0\) for some \(h\).  With \(\mathbf c=\beta\mathbf 1\) and
\(\beta>\|\mathbf d\|_\infty\), the vectors \(\mathbf c\) and \(\mathbf c+\mathbf d\)
are nonnegative, have equal projected observations, and have different signals in
block \(h\).  The rank form holds because \(\dim\sum_gV_g=\operatorname{rank}(A)\),
with equality to \(\sum_g\dim V_g\) exactly when the sum is direct.

\emph{(b, c)}  Let \(B\) index a basis of \(U\) chosen from the columns, write
\(\mathbf a_e=\sum_{b\in B}T_{be}\mathbf a_b\) for \(e\notin B\), and let
\(S_1,\dots,S_m\) be the vertex sets of the connected components of the graph that
joins \(e\notin B\) to every \(b\) with \(T_{be}\ne0\) (isolated vertices are
components).

\emph{Step 1: \(\mathcal S=\{S_i\}\) is identifiable.}  Let \(B_i=B\cap S_i\).  Every
\(e\in S_i\setminus B\) has all basis columns with nonzero coefficient in \(S_i\),
so \(\mathbf a_e\in\operatorname{span}\{\mathbf a_b:b\in B_i\}\) and
\(\operatorname{rank}(A_{S_i})=|B_i|\).  Hence
\(\sum_i\operatorname{rank}(A_{S_i})=|B|=\operatorname{rank}(A)\), and (a) applies.

\emph{Step 2: every identifiable \(\mathcal P\) coarsens \(\mathcal S\).}  Suppose
some \(S_i\) meets two blocks of \(\mathcal P\).  Since \(S_i\) is connected, some
edge \(\{e,b\}\) with \(T_{be}\ne0\) has \(e\in g\) and \(b\in h\) for distinct blocks
\(g,h\).  Let \(d_e=1\), \(d_{b'}=-T_{b'e}\) for \(b'\in B\), and \(d_j=0\)
otherwise; then \(A\mathbf d=0\), and
\(A_h\mathbf d_h=-\sum_{b'\in B\cap h}T_{b'e}\mathbf a_{b'}\) is a combination of
linearly independent columns with the nonzero coefficient \(-T_{be}\), so
\(A_h\mathbf d_h\ne0\).  The interior-point construction contradicts
identifiability.

\emph{Step 3: coarsening preserves identifiability.}  Grouping the summands of a
direct sum gives a direct sum, and \(V_{g\cup h}=V_g+V_h\).

Steps 1--3 show that the identifiable partitions are exactly the coarsenings of
\(\mathcal S\), so \(\mathcal S\) is the unique finest identifiable partition and
does not depend on \(B\).  This proves (b) and the graph characterization in (c).
That the components of this graph are the connected components of the column
matroid is \citet{krogdahl1977dependence}; see also \citet{oxley2011matroid}.

\emph{(d)}  A partition whose blocks are unions of blocks of \(\mathcal D\) is
identifiable exactly when it is a common coarsening of \(\mathcal P^\star\) and
\(\mathcal D\).  The finest common coarsening \(\mathcal P^\star\vee\mathcal D\) is
one, by Step 3.  \(\square\)

\begin{proposition}[Sums of coefficients]
\label{prop:sum_identifiable}
Call \(\mathcal P\) \emph{sum-identifiable} if \(A\mathbf c_1=A\mathbf c_2\) with
\(\mathbf c_1,\mathbf c_2\in\mathbb R_+^J\) implies
\(\sum_{j\in g}c_{1,j}=\sum_{j\in g}c_{2,j}\) for every block \(g\).  This holds if and
only if every block indicator lies in \(\operatorname{row}(A)\); minimal
sum-identifiable partitions need not be unique, and for rational \(A\) deciding whether
one with at least two blocks exists is NP-complete.
\end{proposition}

\noindent\textbf{Proof of Proposition~\ref{prop:sum_identifiable}.}
\(\mathcal P\) is sum-identifiable when \(\mathbf 1_g^\top\mathbf c\) is determined by
\(A\mathbf c\) for every block, on \(\mathbb R_+^J\).  By the interior-point construction, a linear functional
\(\boldsymbol\lambda^\top\mathbf c\) is determined exactly when
\(\boldsymbol\lambda\perp\ker(A)\), that is
\(\boldsymbol\lambda\in\operatorname{row}(A)\), so \(\mathcal P\) is sum-identifiable
exactly when every block indicator lies in \(\operatorname{row}(A)\).

\emph{Non-uniqueness.}  Let \(A\) have rows \((1,1,0,0)\), \((0,0,1,1)\) and
\((1,0,1,0)\).  Its row space contains \((0,1,0,1)\) but no unit vector and not
\((1,0,0,1)\), so the sum-identifiable partitions are \(\{E\}\),
\(\{\{1,2\},\{3,4\}\}\) and \(\{\{1,3\},\{2,4\}\}\); the last two are both minimal.
For the signal target on the same \(A\), \(\ker(A)\) is spanned by
\((1,-1,-1,1)\), the four columns form one circuit, and
\(\mathcal P^\star=\{E\}\).

\emph{NP-completeness.}  The problem ``given rational \(A\), is there a
sum-identifiable partition with at least two blocks?'' is in NP, since
a partition can be checked with rank computations.  For hardness, take an instance
\(w_1,\dots,w_n\) of \textsc{Partition} \citep{garey1979computers} with
\(\sum_iw_i=2T\), let \(\mathbf d=(w_1,\dots,w_n,-T,-T)\in\mathbb Z^{n+2}\), and let
\(A\) be a rational \((n+1)\times(n+2)\) matrix whose rows span
\(\mathbf d^\perp\), so that \(\operatorname{row}(A)=\mathbf d^\perp\).  Because
\(\mathbf 1_E^\top\mathbf d=0\), a partition with at least two blocks exists exactly
when some nonempty proper \(g\subset E\) has \(\mathbf 1_g^\top\mathbf d=0\) (then
\(\{g,E\setminus g\}\) qualifies).  If \(g\) contains neither of the last two
coordinates, \(\mathbf 1_g^\top\mathbf d>0\); if it contains both,
\(\mathbf 1_g^\top\mathbf d=0\) forces \(g=E\); if it contains exactly one, the
condition is \(\sum_{i\in g,\,i\le n}w_i=T\), a solution of \textsc{Partition}.
\(\square\)

\subsection{Structure of Reportable Partitions}

\begin{proposition}[Structure of reportable partitions]
\label{prop:reportable_structure}
\leavevmode
\begin{enumerate}[label=(\alph*),leftmargin=*,itemsep=0pt,topsep=2pt]
  \item The one-block partition \(\{E\}\) is reportable at every \(\tau_\theta\).
  \item If \(g\) and \(h\) are blocks of an identifiable partition, then
  \(1/s_{g\cup h}\le1/s_g+1/s_h\), and reportable partitions are not closed under
  coarsening.
  \item Minimal reportable partitions need not be unique.
\end{enumerate}
\end{proposition}

\begin{remark}[Implication]
Proposition~\ref{prop:reportable_structure} needs only Assumption~\ref{ass:model}.  By
(a) a reportable partition always exists: in the worst case IASA reports the total
signal of all sources.  By (b), merging two blocks does not always help: the merged
block can be closer to the remaining sources than either part, so IASA checks every
union of atoms rather than merging the least separated pair.  By (c), more than one
partition can have the most blocks; IASA therefore states its tie-break (the largest
smallest separation) and lists the remaining ties.
\end{remark}

\subsection{An Oblique Projection Lemma}

\begin{lemma}
\label{lem:oblique}
Let \(U=V\oplus W\) with \(V\ne\{0\}\ne W\), let \(\theta\in(0,\pi/2]\) be the smallest
principal angle between \(V\) and \(W\), and let \(\pi\) be the projection onto
\(V\) along \(W\), defined on \(U\).  Then \(\|\pi\|_2=1/\sin\theta\).
\end{lemma}

\begin{proof}
Write \(\mathbf u=\mathbf x+\mathbf w\) with \(\mathbf x\in V\), \(\mathbf w\in W\).
By the definition of \(\theta\),
\(|\mathbf x^\top\mathbf w|\le\cos\theta\,\|\mathbf x\|_2\|\mathbf w\|_2\), so
\[
\|\mathbf u\|_2^2\ge\|\mathbf x\|_2^2+\|\mathbf w\|_2^2
-2\cos\theta\|\mathbf x\|_2\|\mathbf w\|_2
=(\|\mathbf w\|_2-\cos\theta\|\mathbf x\|_2)^2+\sin^2\theta\|\mathbf x\|_2^2
\ge\sin^2\theta\|\mathbf x\|_2^2 ,
\]
which gives \(\|\pi\mathbf u\|_2\le\|\mathbf u\|_2/\sin\theta\).  For unit principal
vectors \(\mathbf x_0\in V\), \(\mathbf w_0\in W\) with
\(\mathbf x_0^\top\mathbf w_0=\cos\theta\), the vector
\(\mathbf u=\mathbf x_0-\cos\theta\,\mathbf w_0\) has \(\|\mathbf u\|_2=\sin\theta\)
and \(\pi\mathbf u=\mathbf x_0\), so the bound is attained
\citep{ipsen1995angle}.
\end{proof}

If \(V=\{0\}\), then \(\pi=0\); if \(W=\{0\}\), \(\pi\) is the identity on \(U\).  Both
cases are consistent with the convention \(s_g=1\).

\subsection{Group-Signal Error, Reportable Partitions and Stability}

\noindent\textbf{Proof of Theorem~\ref{thm:group_error}.}
\emph{(a)}  By Theorem~\ref{thm:groupings}(a), \(U=\bigoplus_{g}V_g\), and
\(W_g=\bigoplus_{h\ne g}V_h\).  Let \(\pi_g\) be the projection of \(U\) onto
\(V_g\) along \(W_g\).  For any \(\mathbf z\),
\(A\mathbf z=\sum_gA_g\mathbf z_g\) with \(A_g\mathbf z_g\in V_g\), so
\(A_g\mathbf z_g=\pi_g(A\mathbf z)\).  All least-squares solutions have the fitted
signal \(P_A\yt\), and all NNLS solutions have the fitted signal \(\Pi_C(\yt)\),
which is unique because \(C\) is closed and convex; hence the group signals do not
depend on the solution.  Moreover
\(A_g(\widehat{\mathbf c}_g-\mathbf c_g)=\pi_g\bigl(A(\widehat{\mathbf c}-\mathbf c)\bigr)\),
where \(A(\widehat{\mathbf c}-\mathbf c)=P_A\widetilde{\boldsymbol\epsilon}\) for
least squares and \(\|A(\widehat{\mathbf c}-\mathbf c)\|_2\le
\|P_A\widetilde{\boldsymbol\epsilon}\|_2\) for NNLS (proof of
Proposition~\ref{prop:noise_robust}).  Lemma~\ref{lem:oblique} with \(V=V_g\),
\(W=W_g\) gives \(\|\pi_g\|_2\le1/s_g\), with equality when \(V_g\) and \(W_g\) are
nonzero, which proves~\eqref{eq:group_bound}.

\emph{(b)}  Since \(U\subseteq\operatorname{col}(Q)^\perp\),
\(P_A\widetilde{\boldsymbol\epsilon}=P_AP_Q^\perp\boldsymbol\epsilon
=P_A\boldsymbol\epsilon\).  The least-squares error of block \(g\) is
\(M\boldsymbol\epsilon\) with \(M=\pi_gP_A\), a matrix of rank \(r_g\) with
\(\|M\|_2=\|\pi_g\|_2=1/s_g\) when \(r_g\ge1\); when \(r_g=0\), \(M=0\) and the
error is zero.  Hence
\(\mathbb E\|M\boldsymbol\epsilon\|_2^2=\sigma^2\|M\|_F^2\le\sigma^2r_g\|M\|_2^2
=\sigma^2r_g/s_g^2\), with equality when \(r_g=1\) because \(M\) then has one nonzero
singular value.  The map \(\boldsymbol\epsilon\mapsto\|M\boldsymbol\epsilon\|_2\) is
\(\|M\|_2\)-Lipschitz, so the Gaussian concentration inequality of
Tsirelson, Ibragimov and Sudakov \citep{boucheron2013concentration} gives
\(\|M\boldsymbol\epsilon\|_2\le\mathbb E\|M\boldsymbol\epsilon\|_2
+\sigma\|M\|_2\sqrt{2\log(1/\delta)}\) with probability at least \(1-\delta\), and
\(\mathbb E\|M\boldsymbol\epsilon\|_2\le\sigma\|M\|_F\le\sigma\sqrt{r_g}/s_g\); this gives a
high-probability least-squares bound with \(\sqrt{r_g}\) in place of
\(\sqrt{\operatorname{rank}(A)}\), which the theorem does not state.

\emph{(c)}  As in (b), \(P_A\widetilde{\boldsymbol\epsilon}=P_A\boldsymbol\epsilon\).
Write \(\boldsymbol\epsilon=\Sigma^{1/2}\mathbf z\) with \(\mathbf z\sim N(0,I_N)\)
(\(\Sigma=\sigma^2I\) under Assumption~\ref{ass:noise}) and
\(f(\mathbf z)=\|P_A\Sigma^{1/2}\mathbf z\|_2\).  The map \(f\) is Lipschitz with
constant \(\|P_A\Sigma^{1/2}\|_2\le\|\Sigma\|_2^{1/2}\), and
\(\mathbb Ef\le(\mathbb Ef^2)^{1/2}=\|P_A\Sigma^{1/2}\|_F
\le\sqrt{\operatorname{rank}(A)}\,\|\Sigma\|_2^{1/2}\), because
\(P_A\Sigma^{1/2}\) has rank at most \(\operatorname{rank}(A)\).  The Gaussian
concentration inequality \citep{boucheron2013concentration} gives, with probability at
least \(1-\delta\),
\(\|P_A\widetilde{\boldsymbol\epsilon}\|_2\le\|\Sigma\|_2^{1/2}
\bigl(\sqrt{\operatorname{rank}(A)}+\sqrt{2\log(1/\delta)}\bigr)\).  On this event,
part (a), which holds for every noise vector, applies to every block of every
identifiable partition at once.  \(\square\)

\noindent\textbf{Proof of Proposition~\ref{prop:reportable_structure}.}
\emph{(a)}  \(\{E\}\) is identifiable and \(s_E=1\) by convention.
\emph{(b)}  With \(\pi_g\) as above, \(\pi_{g\cup h}=\pi_g+\pi_h\), because both send
\(\sum_k\mathbf x_k\) (with \(\mathbf x_k\in V_k\)) to \(\mathbf x_g+\mathbf x_h\).  By
Lemma~\ref{lem:oblique},
\(1/s_{g\cup h}=\|\pi_g+\pi_h\|_2\le1/s_g+1/s_h\) when \(V_{g\cup h}\ne\{0\}\); when
\(V_{g\cup h}=\{0\}\), \(s_{g\cup h}=1\) and the inequality holds because
\(1/s_g,1/s_h\ge1\).  For the \(4\times4\) operator
\[
\begin{pmatrix}
0.248&-0.142&0.165&0.828\\
0.899&-0.775&-0.917&0.229\\
-0.356&-0.381&0.243&-0.507\\
-0.063&-0.484&-0.270&0.069
\end{pmatrix}
\]
the four singletons have separations of at least \(0.3\), while the blocks of
\(\{\{1,2\},\{3,4\}\}\) have separation below \(0.3\), so at \(\tau_\theta=0.3\) a
reportable partition has a coarsening that is not reportable.
\emph{(c)}  Let \(\at_1=(1,0,0)\), \(\at_2=(\cos a,\sin a,0)\) and
\(\at_3=(\cos a,0,\sin a)\) with \(a=0.3\), and \(\tau_\theta=0.25\).
\(\operatorname{span}(\at_1,\at_2)\) is the plane \(x_3=0\), at angle
\(a\) from \(\at_3\), so both blocks of \(\{\{1,2\},\{3\}\}\) have separation
\(\sin a=0.296\); by symmetry so do those of \(\{\{1,3\},\{2\}\}\).  The angle
between \(\at_1\) and \(\operatorname{span}(\at_2,\at_3)\) has sine
\(\sin a/\sqrt{1+\cos^2a}=0.214\), which is \(s_1\) and also the separation of
\(\{2,3\}\).  At \(\tau_\theta=0.25\) the partitions \(\{\{1,2\},\{3\}\}\) and
\(\{\{1,3\},\{2\}\}\) are reportable, neither refines the other, and the only finer
partition, the three singletons, is not reportable because \(s_1=0.214\).  \(\square\)

\noindent\textbf{Minimal reportable partitions.}  Let the atoms partition \(E\).  A
partition into unions of atoms whose blocks all have \(s_h>0\) is identifiable: \(s_h>0\)
means \(V_h\cap W_h=\{0\}\) with \(W_h=\sum_{k\ne h}V_k\), and a sum of subspaces is direct
exactly when each summand meets the sum of the others only in \(\{0\}\).  So a partition
into unions of atoms is reportable exactly when every block has \(s_h\ge\tau_\theta\).
Since \(s_h\) depends only on \(h\), a reportable partition has a strictly finer reportable
partition into unions of atoms exactly when one of its blocks can be split into unions
of atoms that all have separation at least \(\tau_\theta\).  This gives the
characterization used in step 5 of Algorithm~\ref{alg:iasa}.

\noindent\textbf{Computing minimal reportable partitions.}  With \(q\) atoms, the
separations of all nonempty proper unions of atoms take \(2^{q-1}-1\) principal-angle
computations, because \(s_h=s_{E\setminus h}\) and \(s_E=1\).  From the flags \(s_h\ge\tau_\theta\), dynamic
programming over subsets that contain their lowest atom decides for every union \(S\)
whether it can be partitioned into reportable unions, and hence which reportable unions
cannot be split, in \(O(3^q)\) time.  The same recursion, maximizing first the number
of blocks and then the smallest separation, returns the reported partition in
\(O(3^q)\) time, because \(\min(s_h,\cdot)\) preserves the order of the second criterion.
Listing all minimal reportable partitions, which IASA does to report ties, adds work for
each listed partition.  The greedy rule merges the least separated block with the
partner that maximizes the separation of the merged block; merging \(g\) and \(h\)
leaves every other block's separation unchanged, because \(s_k\) depends only on \(k\),
and the rule stops within \(q-1\) merges because \(s_E=1\).  It has no optimality
guarantee.

\noindent\textbf{Proof of Proposition~\ref{prop:grouping_stability}.}
For unit vectors let \(d(\mathbf x,\mathbf y)=\arccos|\mathbf x^\top\mathbf y|\),
the angle between the lines they span, which is a metric on lines.  Then
\(\theta_g=\min\{d(\mathbf x,\mathbf w):\mathbf x\in V_g,\mathbf w\in W_g\}\).
For a full-column-rank \(M\) and \(\|\Delta\|_2\le\sigma_{\min}(M)/2\), so that the
arcsine arguments below are at most \(1\): a unit
\(\mathbf x'=(M+\Delta)\mathbf z/\|(M+\Delta)\mathbf z\|_2\) in
\(\operatorname{col}(M+\Delta)\) is within distance
\(\|\Delta\mathbf z\|_2/\|(M+\Delta)\mathbf z\|_2\le
\|\Delta\|_2/(\sigma_{\min}(M)-\|\Delta\|_2)\) of \(\operatorname{col}(M)\), and a unit
\(\mathbf x=M\mathbf z/\|M\mathbf z\|_2\) is within distance
\(\|\Delta\|_2/\sigma_{\min}(M)\) of \(\operatorname{col}(M+\Delta)\).  For a unit
vector the distance to a subspace is the sine of its angle to that subspace, so in
both directions the nearest line is within angle
\(\arcsin(\|\Delta\|_2/(\sigma_{\min}(M)-\|\Delta\|_2))\).  Applied to \(M=A_g\) and
\(M=A_{E\setminus g}\) this gives the angles \(\beta_g\) and \(\beta_{E\setminus g}\).  If
\(\mathbf x',\mathbf w'\) attain \(\theta_g(A')\) and \(\mathbf x\in V_g\),
\(\mathbf w\in W_g\) are their nearest lines, the triangle inequality gives
\(\theta_g(A)\le d(\mathbf x,\mathbf w)\le\beta_g+\beta_{E\setminus g}+\theta_g(A')\); the same argument
from \(A\) to \(A'\) gives the reverse inequality.  For the second statement, \(\|\Delta\|_2\le\sigma_J(A)/2\) gives
\(\sigma_J(A')\ge\sigma_J(A)/2>0\), so both operators have full column rank, every
partition is identifiable for both, and both have \(\mathcal P^\star\) equal to the
singletons, so the atoms coincide.  For a nonempty proper union \(g\) of atoms,
\(\sigma_{\min}(A_g)\ge\sigma_J(A)\ge2\|\Delta\|_2\ge2\|\Delta_g\|_2\), and the same for
\(E\setminus g\), so the first statement applies and the margin condition keeps
\(\theta_g\) on the same side of \(\arcsin\tau_\theta\); \(s_E=1\) for both.  By the
characterization of minimal reportable partitions above, the reportable and the minimal
reportable partitions coincide.  \(\square\)
